%\documentclass[12pt,oneside]{amsart}
\documentclass{scrartcl}

%\documentclass{article}

\usepackage{amsthm,amsmath,amssymb}
\usepackage[numbers]{natbib}

%\usepackage[colorlinks]{hyperref}
\usepackage{hyperref}
\hypersetup{
    colorlinks,%
    citecolor=black,%
    filecolor=black,%
    linkcolor=black,%
    urlcolor=black
}
\usepackage{hypernat}
\usepackage[top=2.5cm, bottom=2.5cm, left=2.8cm,
right=2.8cm]{geometry}
\usepackage{graphicx}


%\setlength{\parskip}{0pt}
%\setlength{\parsep}{0pt}
%\setlength{\headsep}{0pt}
%\setlength{\topskip}{0pt}
%\setlength{\topmargin}{0pt}
%\setlength{\topsep}{0pt}
%\setlength{\partopsep}{0pt}

\linespread{1}

%\usepackage{marvosym}

%\textwidth = 6.5 in \textheight = 9.2 in \oddsidemargin = 0.0 in
%\evensidemargin = 0.0 in \topmargin = -0.0 in \headheight = 0.0 in
%\headsep = 0.0in \parskip = 0.0in \parindent = 0.0in
%\def\baselinestretch{0.871}


\newtheorem{theorem}{Theorem}[section]
\newtheorem{proposition}[theorem]{Proposition}
\newtheorem{problem}[theorem]{Problem}
\newtheorem{fact}[theorem]{Fact}
\newtheorem{lemma}[theorem]{Lemma}
\newtheorem{defn}[theorem]{Definition}
\newtheorem{corollary}[theorem]{Corollary}
\newtheorem{remark}{Remark}[section]
\newtheorem{example}[theorem]{Example}

\newcommand{\E}{{\mathbb E}}
\newcommand{\se}{{\mathcal{E}}}
\newcommand{\R}{{\mathbb R}}
\renewcommand{\P}{{\mathbb P}}
\newcommand{\ac}{{\mathcal{A}}}
\newcommand{\A}{{\mathcal{A}}}
\newcommand{\B}{{\mathcal{B}}}
\newcommand{\C}{{\mathcal{C}}}
\newcommand{\M}{{\mathcal{M}}}
\newcommand{\Mf}{{\mathfrak{M}}}
\newcommand{\T}{{\mathcal{T}}}
\newcommand{\Z}{{\mathcal{Z}}}
\newcommand{\K}{{\mathcal{K}}}
\newcommand{\lc}{{\mathcal{L}}}
\newcommand{\Ps}{{\mathcal{P}}}
\newcommand{\G}{{\mathcal{G}}}
\newcommand{\pa}{{\mathring{p}}}
\newcommand{\F}{{\cal F}}
\newcommand{\samp}{{\mathcal{X}}}
\newcommand{\X}{{\mathcal{X}}}
\newcommand{\hi}{{\hat{\Phi}}}
\newcommand{\data}{{\text{data}}}
\newcommand{\relu}{{\text{ReLU}}}

\newcommand{\ridge}{{\hat{\beta}_{\text{ridge}}(\lambda)}}
\newcommand{\lasso}{{\hat{\beta}_{\text{lasso}}(\lambda)}}
\newcommand{\ridgemu}{{\hat{\mu}_t^{\text{ridge}}(\lambda)}}
\newcommand{\lassomu}{{\hat{\mu}_t^{\text{lasso}}(\lambda)}}


\newcounter{rcnt}[section]
\renewcommand{\thercnt}{(\roman{rcnt})}

\def\qt#1{\qquad\text{#1}}

\def\argmin{\mathop{\rm argmin}}
\def\argmax{\mathop{\rm argmax}}


\setlength{\parskip}{1.4 \medskipamount}
%\sloppy
%\linespread{1.3}

\begin{document}

\title{STAT 238 - Bayesian Statistics  \\
  Lecture Twenty Six} 
\subtitle{Spring 2026, UC Berkeley}

\author{Aditya Guntuboyina}


\date{03 April 2026}

\maketitle

\section{The Gibbs Sampler}
Let us introduce some notation. The goal is to obtain
samples from a target distribution with density $\pi$ (in our
applications, $\pi$ will be the posterior density). $\pi$ represents
the density on $\R^d$ of a random vector $\Theta$ of size $d \times
1$:  
\begin{equation*}
   \Theta \sim \pi. 
 \end{equation*}
 We assume that $\Theta$ can be decomposed as:
 \begin{equation*}
   \Theta = (\Theta_{(1)}, \dots, \Theta_{(k)}) 
 \end{equation*}
 for $k$ subvectors $\Theta_{(1)}, \dots, \Theta_{(k)}$. The $j^{th}$
 subvector is $\Theta_{(j)}$. We shall also denote the vector obtained
 by dropping the $j^{th}$ subvector $\Theta_{(j)}$ from $\Theta$ by
 $\Theta_{-(j)}$. For example, 
 \begin{equation*}
   \Theta_{-(1)} = (\Theta_{(2)}, \dots, \Theta_{(k)}) ~~ \text{  and
   } ~~ \Theta_{-(2)} = (\Theta_{(1)}, \Theta_{(3)}, \dots,
   \Theta_{(k)}). 
 \end{equation*}
 Gibbs sampling works according to the following algorithm.
 \begin{enumerate}
 \item Initialize with an arbitrary value $\theta^{(0)}$.
 \item Repeat the following for $t = 1, \dots, N$ (for a large number
   $N$):
   \begin{enumerate}
   \item Pick $J$ uniformly at random from $1, \dots, k$. Suppose $J$
     turned out to be $j$. 
   \item  Set $\theta^{(t+1)}_{-(j)} = \theta^{(t)}_{-(j)}$ i.e.,
     $\theta^{(t+1)}_{(i)} = \theta^{(t)}_{(i)}$ for $i \neq j$.
    \item $\theta^{(t+1)}_{(j)}$ is randomly drawn according to the
      conditional distribution of $\Theta_{(j)}$  given $\Theta_{-(j)}
      = \theta^{(t)}_{-(j)}$. 
   \end{enumerate}
 \end{enumerate}
This algorithm assumes that we have some way of generating
observations from the conditional distribution of $\Theta_{(j)}$ given
$\Theta_{-(j)}$ for each $j$.

\section{Examples of the Gibbs Sampler}

In the first two examples below, the posterior can be written in
closed form so MCMC is not really necessary. These are still useful
for evaluating the performance of the Gibbs sampler. The third
example is more non-trivial and illustrates the usefulness of the
Gibbs sampler.    

\subsection{Bivariate Normal}
We want to sample $(\theta_1, \theta_2)$ from the following bivariate
distribution: 
  \begin{equation*}
    \begin{pmatrix} \theta_1 \\ \theta_2 \end{pmatrix}
    \Biggr| \begin{pmatrix}y_1 \\ y_2 \end{pmatrix} \sim N
    \left(\begin{pmatrix}y_1 \\ y_2 \end{pmatrix}, \begin{pmatrix} 1 & \rho \\ \rho &
    1 \end{pmatrix} \right). 
\end{equation*}
Here $y = (y_1, y_2)^T$ is fixed. One can, of course, directly sample from
this normal distribution without any MCMC. We can use this example to
illustrate the performance of the Gibbs sampler. To use the Gibbs
sampler, observe that  
  \begin{equation*}
    \theta_1 | \theta_2, y \sim N(y_1 + \rho(\theta_2 - y_2), 1 -
    \rho^2) ~~ \text{ and } ~~ \theta_2 | \theta_1, y \sim N(y_2 + \rho(\theta_1 - y_1), 1 -
    \rho^2)
  \end{equation*}
  The Gibbs sampler algorithm then would be the following:
  \begin{enumerate}
  \item Initialize at arbitrary $\theta_1^{(0)}$ and
    $\theta_2^{(0)}$.
  \item Repeat the following for $t = 1, \dots, N$: 
    \begin{enumerate}
    \item Pick $J$ to be either 1 or 2 with probability $0.5$.
     \item If $J = 1$, take $\theta_2^{(t+1)} = \theta_2^{(t)}$ and
       generate $\theta_1^{(t+1)}$ from the normal distribution with
       mean $y_1 + \rho (\theta_2^{(t)} - y_2)$ and variance $1 -
       \rho^2$.
      \item   If $J = 2$, take $\theta_1^{(t+1)} = \theta_1^{(t)}$ and
       generate $\theta_2^{(t+1)}$ from the normal distribution with
       mean $y_2 + \rho (\theta_1^{(t)} - y_1)$ and variance $1 -
       \rho^2$.
    \end{enumerate}
  \end{enumerate}
  This example is also useful for illustrating one potential problem
  with the Gibbs sampler. Suppose $\rho = 1$. In this case $\theta_1 -
  y_1 \sim N(0, 1)$ and $\theta_2 - y_2 = \theta_1 - y_1$. The Gibbs
  iterates in this case will be
  \begin{equation*}
    \theta_1^{(t+1)} - y_1 = \theta_2^{(t)} - y_2 ~~~ \text{if $J = 1$}
  \end{equation*}
  and
  \begin{equation*}
    \theta_2^{(t+1)} - y_2 = \theta_1^{(t)} - y_1 ~~~ \text{if $J = 2$}
  \end{equation*}
  This results in the chain not moving after the first iterate. For
  example, if the first step is: 
  \begin{equation*}
    \begin{pmatrix} \theta_1^{(0)} \\ \theta_2^{(0)}\end{pmatrix}
    \rightarrow   \begin{pmatrix} y_1 - y_2 + \theta_2^{(0)}\\
      \theta_2^{(0)}\end{pmatrix},  
  \end{equation*}
  then the chain will stay at $\begin{pmatrix} y_1 - y_2 + \theta_2^{(0)}\\
      \theta_2^{(0)}\end{pmatrix}$ forever. On the other hand, if the
    first step is
  \begin{equation*}
    \begin{pmatrix} \theta_1^{(0)} \\ \theta_2^{(0)}\end{pmatrix}
    \rightarrow   \begin{pmatrix} \theta_1^{(0)}\\
      y_2 - y_1 + \theta_1^{(0)}\end{pmatrix},  
  \end{equation*}
  then the chain will stay forever at  $\begin{pmatrix} \theta_1^{(0)}\\
      y_2 - y_1 + \theta_1^{(0)}\end{pmatrix}$. This chain will obviously
    not converge unless the initial iterate already has the right
    distribution. A similar problem arises when $\rho = -1$. 

    The lesson to be drawn from this is that the Gibbs sampler has
    issues of convergence when there is high correlation between the
    variables. 

\subsection{Linear Regression}
  Consider the linear regression model that we discussed
  previously. We observe data $(y_1, x_1), \dots, (y_n, x_n)$ with
  $y_1, \dots, 
  y_n$ real-valued and $x_1, \dots, x_n \in \R^p$. We treat
  $x_1, \dots, x_n$ as non-random and $y_1, \dots, y_n$ as random. The
  parameter is $\theta \in \R^{p+1}$ with $\theta = (\beta, \sigma)$
  (here $\beta \in \R^p$ represents the vector of regression coefficients and
  $\sigma$ is the noise standard deviation). Consider the model 
  \begin{equation*}
   y_i  \sim N(x_i' \beta, \sigma^2)
 \end{equation*}
 with $y_1, \dots, y_n$ being independent conditional on $\theta$
 along with the (improper) prior:
 \begin{equation*}
  \pi(\beta, \sigma) \propto \frac{1}{\sigma} I\{\sigma > 0\}. 
\end{equation*}
The exact joint posterior is given by
\begin{equation}\label{expo_reg}
  \pi(\beta, \sigma | \text{data})  \propto
  \sigma^{-(n+1)} \exp \left(\frac{-1}{2 \sigma^2} \|Y  - X \beta\|^2
  \right) I\{\sigma > 0\}
\end{equation}
We can directly sample from this posterior (e.g., by first sampling
from the marginal posterior of $\sigma$ and then from the conditional
posterior of $\beta$ given $\sigma$. Alternatively, one can first
sample from the marginal posterior of $\beta$ and then from the
conditional posterior of $\sigma$ given $\beta$).

Now let us use the Gibbs sample to simulate from
\eqref{expo_reg}. Note that:
   \begin{equation*}
     \pi(\beta | \sigma, \text{data}) \propto \exp \left(\frac{-1}{2
         \sigma^2} \|Y - X \beta\|^2 \right) \propto \exp \left(\frac{-1}{2
         \sigma^2} \|X \beta - X \hat \beta\|^2 \right) = \exp \left(\frac{-1}{2
         \sigma^2} (\beta - \hat \beta)' (X' X) (\beta - \hat \beta) \right)
   \end{equation*}
   which gives
   \begin{equation*}
     \beta | \sigma, \text{data} \sim N_p \left(\hat \beta, \sigma^2
       (X' X)^{-1} \right). 
   \end{equation*}
Here $\hat{\beta}$ is the least squares estimator. For the conditional
posterior of $\sigma$ given $\beta$, note that
\begin{equation*}
  \pi(\sigma | \beta, \text{data}) \propto   \sigma^{-(n+1)} \exp
  \left(\frac{-1}{2 \sigma^2} \|Y  - X \beta\|^2 \right) I\{\sigma >
  0\}. 
\end{equation*}
Using the formula $f_{1/\sigma^2}(x) = f_{\sigma}(x^{-1/2}) x^{-3/2}$,
it is now easy to check that the density of $1/\sigma^2$  conditional
on $\beta$ and the data is proportional to  
\begin{equation*}
    x^{(n-2)/2}  \exp \left(\frac{-x}{2} \|Y - X \beta\|^2 \right) I
    \{x > 0\}
  \end{equation*}
  which means that
  \begin{equation*}
    \frac{1}{\sigma^2} \biggr| \beta, \text{data} \sim \text{Gamma}
    \left(\frac{n}{2}, \frac{1}{2} \|Y - X\beta\|^2 \right). 
  \end{equation*}
   The Gibbs sampler algorithm is thus: 
   \begin{enumerate}
   \item Initialize at arbitrary $\beta^{(0)}$ and $\sigma^{(0)}$.
   \item Repeat the following for $t = 1, \dots, N$:
     \begin{enumerate}
     \item Pick $J$ to be either $1$ or $2$ with equal probability.
     \item If $J =1$, take $\beta^{(t+1)} \sim N_p(\hat{\beta},
       (\sigma^{(t)})^2 (X' X)^{-1})$ and $\sigma^{(t+1)} =
       \sigma^{(t)}$.
     \item If $J = 2$, take $\beta^{(t+1)} = \beta^{(t)}$. Generate an observation
     $x$ from the Gamma distribution with parameters $n/2$ and
     $\|Y - X\beta^{(t)}\|^2/2$. Then take $\sigma^{(t+1)} =
     x^{-1/2}$. 
     \end{enumerate}
   \end{enumerate}


\subsection{Probit Regression}
Our next application of the Gibbs sampler is to Probit
Regression which was introduced in Homework Four. The model is given by
the following. We observe data $(y_1, x_1), \dots, (y_n, x_n)$ with 
  $y_1, \dots, y_n$ binary and $x_1, \dots, x_n \in \R^p$. $x_1,
  \dots, x_n$ are treated non-random and $y_1, \dots, y_n$ random. The 
  parameter is $\beta \in \R^p$. The probit regression model assumes
  that $y_1 \dots, y_n$ are independent given $\beta$ with
  \begin{equation*}
    y_i | \beta \sim \text{Bernoulli} \left(\Phi(x_i' \beta) \right)
  \end{equation*}
  where $\Phi(\cdot)$ is the standard normal cdf. Let $\pi(\cdot)$ be
  our prior for $\beta$. The posterior is then given by:
  \begin{equation*}
    \pi(\beta | \text{data}) \propto \pi(\beta) \prod_{i=1}^n
    \left(\Phi(x_i' \beta) \right)^{y_i} \left(1 - \Phi(x_i' \beta)
    \right)^{1-y_i}. 
  \end{equation*}
  Suppose we take the (improper) uniform prior for $\beta$ so that
  the posterior becomes
  \begin{equation*}
    \pi(\beta | \text{data}) \propto \prod_{i=1}^n
    \left(\Phi(x_i' \beta) \right)^{y_i} \left(1 - \Phi(x_i' \beta)
    \right)^{1-y_i}. 
  \end{equation*}
  It is not actually clear how to use the Gibbs sampler on the above
  posterior. We can try decomposing $\beta$ as $(\beta_1, \dots,
  \beta_p)$ and then attempt to obtain samples from the conditional:
  \begin{equation*}
    \pi(\beta_j | \beta_{-j}, \text{data}) \propto \prod_{i=1}^n
    \left(\Phi(x_i' \beta) \right)^{y_i} \left(1 - \Phi(x_i' \beta)
    \right)^{1-y_i}. 
  \end{equation*}
  I am not sure how to obtain samples from the above conditional.

  There is a trick that people use to run the Gibbs sampler for probit
  regression. Let us introduce random variables $w_1, \dots, w_n$
  which, conditional on $\beta$, are independent with
  \begin{equation*}
   w_i | \beta \sim N(x_i' \beta, 1). 
  \end{equation*}
  We then let
  \begin{equation*}
    y_i = I\{w_i > 0\}. 
  \end{equation*}
  It is then clear that $y_1, \dots, y_n$ defined thus have the same
  likelihood as in the probit regression model. Let $w = (w_1, \dots,
  w_n)$. The trick is to use the Gibbs sampler with $(w, \beta)$. We shall see how this is done in the next class. 




%\begin{thebibliography}{99}

%\bibitem[Roberts and Rosenthal(2004)]{RobertsRosenthal2004}
%Roberts, G. O. and Rosenthal, J. S. (2004).
%General state space Markov chains and MCMC algorithms.
%\textit{Probability Surveys}, \textbf{1}, 20--71.

%\bibitem[Meyn and Tweedie(2009)]{MeynTweedie2009}
%Meyn, S. P. and Tweedie, R. L. (2009).
%\textit{Markov Chains and Stochastic Stability}, 2nd ed.
%Cambridge University Press.


  
%\bibitem[Chib and Greenberg(1995)]{chib1995}
%Chib, S. and Greenberg, E. (1995).
%Understanding the Metropolis--Hastings algorithm.
%\textit{The American Statistician}, \textbf{49}, 327--335.


%\bibitem[MacKay and Peto(1995)]{MacKay1995HierarchicalDirichlet}
%D.~J.~C. MacKay and L.~C. Peto,
%\newblock ``A hierarchical Dirichlet language model,''
%\newblock \emph{Natural Language Engineering}, vol.~1,
%no.~3, pp.~289--308, 1995.

%\bibitem[Rubin(1981)]{Rubin1981BayesianBootstrap}
%D.~B. Rubin,
%\newblock ``The Bayesian bootstrap,''
%\newblock \emph{The Annals of Statistics}, vol.~9,
%no.~1, pp.~130--134, 1981.
  
%\bibitem[Fisher(1934)]{Fisher1934}
%R.~A. Fisher,
%\newblock ``Probability, likelihood and quantity of information in the logic of
%uncertain inference,''
%\newblock \emph{Proceedings of the Royal Society of London. Series~A,
%Containing Papers of a Mathematical and Physical Character}, vol.~146,
%no.~856, pp.~1--8, 1934.

%\bibitem[Efron(2010)]{Efron}
%Efron, B. (2010)
%\textit{Large-scale inference: empirical Bayes methods for estimation,
%testing, and prediction}. Cambridge University Press, 2010.

%\bibitem[Shumway and Stoffer(2017)]{ShumwayStoffer}
%Shumway, R. H. and Stoffer, D. S. (2017)
%\textit{Time Series Analysis and Its Applications: With R
%  Examples}. Springer, 4th edition. 
  
%     \bibitem[Jaynes(2003)]{Jaynes} Jaynes, E. T, (2003)
%  \textit{Probability theory: the logic of science}. Cambridge
%  University Press, 2003.
 
%  \bibitem[Sinai(1992)]{Sinai} Sinai, Y. G, (1992)
%  \textit{Probability theory: an introductory course}. Springer
%  Verlag, 1992.  

%\bibitem[{deGroot(1986)}]{DeGrootBlackwell}DeGroot, M. H. (1986)
%       A conversation with David Blackwell.
%\newblock \emph{Statistical Science}, \textbf{1}, 40--53.

%         \bibitem[Mosteller(1987)]{Mosteller} Mosteller, F, (1987)
%  \textit{Fifty challenging problems in probability with
%    solutions}. Courier Corporation, 1987.

%    \bibitem[MacKay(2003)]{MacKay} MacKay, D. J. C., (2003)
%  \textit{Information theory, inference, and learning
%  algorithms}. Cambridge University Press, 2003.

%\bibitem[{Lindley(2013)}]{Lindley}Lindley, D. V. (2013)
%   \textit{Understanding Uncertainty}. John Wiley and sons, 2013.


%\end{thebibliography}







\end{document}

%%% Local Variables:
%%% mode: latex
%%% TeX-master: t
%%% End:
